\section{Finite palette and guard certificates}
\label{app:palette}

This appendix identifies the finite objects used in \cref{lem:palette}.  The
underlying active graphs are four explicit constructors in the source artifact:
\path{state_0m1}, \path{state_00m11}, \path{state_01m10}, and
\path{state_00m10m11}.  They accompany Rudolph's gain/loss
construction~\cite{rudolph2026universal}.
The supplementary certificate archive records the complete graphs, oriented
integer fillings, kernel bases, modular rank witnesses, and Gram determinants.
No numerical eigenvalue threshold is used.

\begin{table}[ht]
\centering
\begin{tabularx}{\linewidth}{@{}>{\raggedright\arraybackslash}Xrrrr>{\raggedright\arraybackslash}X@{}}
\toprule
Active vector & degree & vertices & \(\dim V\) & \(\dim Q\) & certificate mod \(p\) \\
\midrule
\(\ket0-\ket1\) & 1 & 16 & 2 & 8 &
\(\det G=512\) \\
\(\ket{00}-\ket{11}\) & 3 & 33 & 4 & 112 &
\(\det G=655970\) \\
\(\ket{01}-\ket{10}\) & 3 & 33 & 4 & 112 &
signed relabeling \\
\(\ket{00}-\ket{10}-\ket{11}\) & 3 & 41 & 4 & 176 &
\(\det G=797443\) \\
three remaining gain/loss vectors & 3 & 41 & 4 & 176 &
signed relabelings \\
\bottomrule
\end{tabularx}
\caption{Exact active-palette data.  In each determinant row,
\(p=1000003\) and \(G=T_{a,1}^*T_{a,1}\) on \(Q_a\); a nonzero residue
proves injectivity over \(\mathbb Q\).}
\label{tab:palette-certificates}
\end{table}

For the one-qubit atom, exact boundary ranks are \(15,24\) and the Betti
numbers are \((1,1,0)\); a denominator-one 24-term two-chain fills the target
cycle.  At zero weight the target differential pair has rank \(30\), and the
two register cycles plus eight private vectors give its complete
ten-dimensional kernel.  For \(\ket{00}-\ket{11}\), the clique counts in
degrees zero through four are
\((33,241,646,712,272)\), the boundary ranks are
\((32,209,437,272)\), and the Betti numbers are \((1,0,0,3,0)\).
A denominator-one 248-term four-chain fills the target; rank \(596\) of the
zero-weight pair matches the four register and 112 private kernel vectors.
Swapping the second bowtie petals transports this certificate exactly to
\(\ket{01}-\ket{10}\).

For the three-term representative, the clique counts in degrees zero through
five are
\[
(41,322,914,1082,477,30),
\]
and the boundary ranks are \((40,282,632,447,30)\).  Hence the target
degree-three Betti number is three.  A denominator-one 382-term four-chain
fills \(\ket{00}-\ket{10}-\ket{11}\), while three complementary register
cycles remain independent modulo boundaries.  The zero-weight pair has rank
\(902\), matching its kernel \(V\oplus Q\) of dimension \(4+176\).
Signed permutations of bowtie petals transport the graph, filling, kernel,
and private pair to the other three gain/loss vectors.

\subsection{Selected-cycle guards}

We record the guard closure because it is part of the theorem interface, not
an appeal to a general integer-state gadget.  Let \(Y\) implement
\(\pi_\phi\) on register \(R\), let \(B\) be a fresh bowtie with petals
\(S_0,S_1\), and choose \(\beta\in\{0,1\}\).  Add all edges of \(B\), join
every vertex of \(R\) to every vertex of \(B\), and join each private vertex
of \(Y\) only to the selected petal \(S_\beta\).  The resulting flag complex
is
\begin{equation}
\label{eq:guard-union}
Z=(Y*S_\beta)\cup(R*B),\qquad
(Y*S_\beta)\cap(R*B)=R*S_\beta .
\end{equation}
Its new register is \(R'=R*B\).  If \(\gamma_\beta\) is the normalized
oriented cycle of the selected petal, then \(Z\) implements
\[
\pi_\phi\otimes\ket{\beta}\!\bra{\beta},\qquad
Q'=Q\otimes\mathbb C\gamma_\beta .
\]

Indeed, reduced relative chains give
\[
\widetilde C_*(Z,R')
\cong\widetilde C_*(Y,R)\widehat\otimes
\widetilde C_*(S_\beta).
\]
The connecting map therefore tensors the old exact boundary intersection
with \([\gamma_\beta]\), proving exact filling and no additional killed
register class.  At zero private weight the private differential is the
corresponding tensor-sum differential.  Since \(S_\beta\) has reduced
homology only in degree one, the complete target kernel is
\(V'\oplus Q'\); every other bidegree pays a fixed positive guard gap.
Finally, project the old private output pair onto the selected harmonic line.
Both outside differentials vanish under this projection, so
\[
T_Z(\lambda)=\lambda T_{Z,1},\qquad
T_{Z,1}V'=0,\qquad
\|T_{Z,1}(q\otimes\gamma_\beta)\|
\ge b\|q\otimes\gamma_\beta\|.
\]
Removing the unique old private vertex also gives
\(LL^*\preceq N_{\rm priv}\lambda^2I\).  Iterating this construction for the
constant number of required basis guards preserves a fixed finite palette.

The basis projector itself is the special case
\(R\cup(v*S_\beta)\), where a private apex \(v\) cones the selected petal.
Its private zero-weight kernel \(Q\) is zero, so the private-pair condition is
vacuous.  Together with \cref{tab:palette-certificates}, this proves
\cref{lem:palette}.

\subsection{Reproducibility boundary}

The supplementary archive contains the exact witnesses in
\begin{itemize}
\item \path{RUDOLPH_REPRESENTATIVE_BULK_CERTIFICATE.json},
\item \path{REMAINING_ACTIVE_ATOM_CERTIFICATES.json},
\item \path{ACTIVE_HADAMARD_ORBIT_CERTIFICATE.json}, and
\item \path{SELECTED_CYCLE_GUARD_CHECKS.json}.
\end{itemize}
Their checkers recompute the stated ranks, fillings, relabelings, and guard
identities.  These finite checks certify the displayed palette; the analytic
implication from those certificates is the proof of \cref{thm:transfer}.
